\code{LibraryTracking} is used to configure the handling of library events for
its associated Process.

\begin{apient}
static void setDefaultTrackLibraries(bool b)
\end{apient}

\apidesc{
Sets the default handling mechanism for library events across all Process
objects to interrupt-driven (b = true) or polling-driven (b = false).
}

\begin{apient}
static bool getDefaultTrackLibraries()
\end{apient}

\apidesc{
Returns the current default handling mechanism for library events across all
Process objects. A return value of true indicates interrupt-driven while
false indicates polling-driven.
}

\begin{apient}
bool setTrackLibraries(bool b) const
\end{apient}

\apidesc{
Sets the library event handling mechanism for the associated Process object to
interrupt-driven (b = true) or polling-driven (b = false).
Returns true on success and false on failure.
}

\begin{apient}
bool getTrackLibraries() const
\end{apient}

\apidesc{
Returns the current library event handling mechanism for the associated Process
object. A return value of true indicates interrupt-driven while false
indicates polling-driven.
}

\begin{apient}
bool refreshLibraries()
\end{apient}

\apidesc{
Manually polls for queued library events to handle.
Returns true on success and false on failure.
}
